% !TeX root = ../教材题库_全部条目_转写工作稿.tex
% Source: XJ-B1-C06-HW-P257-Y0217
% 第6章 / 6.4 用样本估计总体 / 教材 p.257 / PDF 第264页
% OCR未校对：公式与文字均可能有误；未完成数学TeX转写；禁止直接发给学生
% 原题及相邻原解截图（核对用）：../题目原图/XJ-B1-C06-HW-P257-Y0217-01.webp
% 本文件只在转写完成后作为题干使用；OCR 条目尚不是可靠题干。
\begin{bankitem}{习题6.4 第12题}{第6章 / 6.4 用样本估计总体 / 教材 p.257 / PDF 第264页}{XJ-B1-C06-HW-P257-Y0217}{OCR未校对：公式与文字均可能有误；未完成数学TeX转写；禁止直接发给学生}
\begin{ocrdraft}
[教材p.257 / PDF第264页]
\par
12．在分层抽样时，如果总体分为差冫3，差eZ)层，而且第层抽取的样本
\par
量为屿，第j层的样本均值为，样本方差为寸，丿一1，2，
\par
求证：所有数据的样本均值和方差分别为
\par
1
\par
[屿寸十（
\par
一的21
\end{ocrdraft}
\end{bankitem}
